\chapter{Introduction}
\label{ch:introduction}

The Bitcoin \gls{ln} is a \gls{pcn} built upon the Bitcoin network. 
Since its inception in 2018, the network has experienced remarkable growth, making its evolution a subject 
of significant interest for researchers, protocol engineers, and developers. 
The \gls{ln} can be abstracted as a graph object. 
In April 2026 the public network consists of more than
\num{15000} active nodes and approximately \num{45000} payment channels acting as edges.

Because the \gls{ln} operates as a distributed \gls{p2p} network, 
information regarding its dynamic topology (e.g. the creation of new nodes or channels)
cannot be centrally managed or requested. 
Instead, the network relies on the Lightning Protocol (\gls{bolt}), which defines 
the communication, mainly three types of \term{gossip messages}: \code{channel\_announcement}, \code{node\_announcement} and \code{channel\_update}. 
These messages are continuously broadcast between nodes and form the foundation for network discovery and routing.

Capturing and analysing this distributed data is a technical challenge. 
This thesis introduces \tech{ln-history}, a platform featuring a novel architecture designed 
not only to collect and persist gossip messages (via real-time ingestion and bulk imports) 
but also to enable deep empirical analysis.


\section{Research Contributions}
This thesis presents both the complete system design of \tech{ln-history} and 
an empirical analysis demonstrating the platform's wide range of capabilities.
\begin{itemize}
    \item \textbf{High-Performance Architecture:} 
        The platform features lightning-fast historical snapshot generation (averaging approximately \qty{3}{\second} per snapshot), 
        a capability previously unavailable in the research ecosystem.
    \item \textbf{Real-Time and Historical Analysis:} 
        The platform is utilized to measure real-time metrics such as message volume and propagation lag, 
        while also hosting one of the most complete historical gossip datasets available, 
        making it ideal for tracking the long-term evolution of the \gls{ln}.
    \item \textbf{Open-Source Verifiability:} 
        Ensuring absolute transparency, the entire project, including its source code and infrastructure components, 
        is completely open-source\footnote{\url{https://github.com/ln-history}}.
\end{itemize}


\section{Scope and Conventions}
To ensure clarity and precision throughout this thesis, 
the following technical scopes and typographical conventions are established:

\paragraph{Nomenclature and Units}
Following the convention introduced in~\citetitle[chapter ~1]{antonopoulos_mastering_2021} by~\citeauthor{antonopoulos_mastering_2021}, 
the capitalized word \textit{Bitcoin} refers to the network and technology as a whole, 
whereas the lowercase \textit{bitcoin} refers to the unit of currency (abbreviated as BTC)~\cite{antonopoulos_mastering_2021}. 
One bitcoin is divisible into \num{100000000} \term{satoshis} (sat). 
Because the \gls{ln} supports fractional routing fees, the primary unit of account 
within the protocol is \term{milli-satoshi} (msat). 
In this thesis, all three units (bitcoin, sat, msat) are utilized depending on the relevant scale of the values. 
Additionally, data sizes are denoted using standard abbreviations, 
distinguishing between \textit{bytes} (B) and \textit{bits} (b).

\paragraph{Network Scope}
While it is common to refer to \enquote{the} Lightning Network, 
the protocol does not enforce a single, unified topology. 
The network can technically split up into multiple isolated components. 
Therefore, whenever this thesis refers to the \gls{ln} topology, 
it actually describes the \textit{largest connected component} of the network. 
Additionally, the protocol allows for the creation of \textit{private} (unannounced) channels. 
Because these channels are intentionally hidden from the public gossip protocol, 
they are excluded from the empirical analysis presented in this work.
\newpage


\section*{Thesis Outline}

The structure of the thesis is as follows:

\begin{itemize}
    \item \textbf{Related Work}  
        \hyperref[ch:related-work]{Chapter~\ref*{ch:related-work}} summarizes the academic and developer landscape for Bitcoin and \gls{ln} \gls{p2p} network data analysis. 
        It highlights the various research papers (particularly the \term{time-machine} methodology developed by \citeauthor{decker_lnresearch_2020}~\cite{decker_lnresearch_2020} 
        and utilized by \citeauthor{zabka_empirical_2022}~\cite{zabka_empirical_2022}) that motivated the development of \tech{ln-history}.
    
    \item \textbf{Background} 
        \hyperref[ch:background]{Chapter~\ref*{ch:background}} provides the theoretical foundation of the \gls{ln} protocol, known as \term{BOLT}. 
        It focuses specifically on BOLT\#7, which defines the gossip protocol, 
        while also providing the necessary introductory concepts regarding the Bitcoin system.
    
    \item \textbf{Architecture} 
        \hyperref[ch:architecture]{Chapter~\ref*{ch:architecture}} details the complete design of the \tech{ln-history} platform. 
        It outlines the architectural trade-offs, system requirements, and the database-centric approach
        enabling in-depth analytical queries. 
        The system's components are presented sequentially, matching the data flow: 
        from real-time ingestion via the \service{gossip-publisher-zmq} \gls{cln} plugin, 
        to the \service{gossip-processor}, and finally to the \service{ln-history-database},
        and also details historical bulk imports,
        the \term{RESTful} \service{ln-history-api}, 
        a ready-to-use Python client (\service{ln-history-python-client}), 
        and the platform's observability and deployment strategies.
    
    \item \textbf{Empirical Data Analysis} 
        \hyperref[ch:empirical-data-analysis]{Chapter~\ref*{ch:empirical-data-analysis}} demonstrates the analytical capabilities of the platform through a two-part study. 
        Firstly, it analyses real-time metadata collected over 48 days (January 14th to March 3rd 2026) 
        by two independent collector nodes, measuring propagation lag, message volume, topological overlap, and peer-session durations. 
        Secondly, it performs a historical analysis utilizing remarkable \num{91.3} million gossip messages 
        to construct monthly snapshots from 2018 to April 2026, evaluating channel lifetimes and network centrality metrics.
    
    \item \textbf{Conclusion} 
        \hyperref[ch:conclusion]{Chapter~\ref*{ch:conclusion}} summarizes the empirical findings and discusses their broader implications 
        for the network's efficiency and topology.
    
    \item \textbf{Future Work:} 
        \hyperref[ch:future-work]{Chapter~\ref*{ch:future-work}} outlines how the platform's infrastructure can be leveraged to tackle future research questions. 
        Furthermore, it highlights the latest protocol developments e.g. the Gossip v2 specification (which was just merged on May 4th 2026) 
        and provides an outlook on emerging Layer-2 and Layer-3 scaling solutions within the broader Bitcoin ecosystem.
\end{itemize}
